\documentclass[12pt, draftcls, onecolumn]{IEEEtran}
\usepackage{amsmath}
\usepackage{graphicx}
\usepackage{tabularx} % Put this in your preamble
\usepackage{array}    % For column formatting
\usepackage[utf8]{inputenc}
\usepackage{chemarr}
\usepackage{float}

\usepackage{amssymb}
\usepackage{booktabs}
\usepackage{hyperref}
\usepackage{cleveref}


\title{Virtual Clinical Trial Design for Infective Endocarditis: Evaluating Antibiotic Therapy}
\author{Jacob W. Kammer, Klas Udekwu, and Esteban A. Hernandez-Vargas}
\date{\today} % Use \date{} to leave it blank or specify a custom date

\begin{document}


\maketitle
 


\section{Introduction}
Infective endocarditis is a disease of the endocardial surface of the heart that involves the cardiac valves and the mural endothelium \cite{Li2024}. It is an indolent infection with Staphylococcus aureus (S. aureus) and Viridans streptococci as the most common pathogens \cite{HammondHaley2023}. It is an aggressive infection with an in hospital mortality rate of 15\% to 20\% and a one-year mortality rate of 40\% \cite{Taduru2023}. Complications include congestive heart failure, thromboembolic events, perivalvular and annular abscesses \cite{murdoch2009infective}, sepsis and septic shock. Methicillin-resistant Staphylococcus aureus (MRSA) is the most common cause of bacterial endocarditis in much of the world \cite{HammondHaley2023}.    

Vancomycin has become the cornerstone of antibiotic treatment against MRSA infective endocarditis since the 1970s, and is the first line treatment against MRSA endocarditis \cite{Minter2023}. Vancomycin, a glycopeptide, inhibits cell wall synthesis by preventing polymerization of glycan strands and cross linking between the strands via transpeptidase (PBP) \cite{murray_glycopeptides_2019}. Increased reliance on treatment with vancomycin since the 1980s has led to the development of resistant forms \cite{mcguinness2017vancomycin}. Infections of the heart values have become mixtures of MRSA sensitive and resistant to vancomycin.  Compounding the problem of resistance is the challenge of treating the original source of infection, such as orthopedic hardware, central venous catheters and infection of the bone. It is well understood that S. aureus infections can be chronic and recurrent despite appropriate antibiotic treatment \cite{Kahl2016SCV}. overarching all of these issues is the pivotal role of the immune response in resolving infections. Depending on the clinical scenario treatment failure rates can be as high as 20 to 30 percent.  All of these factors in combination pose a significant clinical challenge that will require an improved understanding.  

Since the end of the 20th century newer antibiotics, such as linezolid have been developed to target the antibiotic resistant S. aureus.  Linezolid, an oxazolidinone antibiotic, has become one of the second line agents used in the treatment of vancomycin resistant MRSA endocarditis \cite{Cong2020}.  Beyond simply adding linezolid to treatment protocols an improved understanding of the factors that lead to clinical resolution (bacteremia resolved) or failure (persistent or relapsing bacteremia) is needed. Understanding how bacterial phenotypic variability, antibiotic combinations, and host factors interact will lead to improved outcomes in bacterial endocarditis.  

Mikkaichi and colleagues \cite{Mikkaichi2019} developed a mathematical model to examine the interplay of phenotypic heterogeneity
with host factors in the persistence of MRSA bacteremia.  Mikkaichi used ensemble modeling and other ML techniques to determine
the key differences between resolving and persistent bacteremia \cite{Mikkaichi2019}.
Two phenotypically distinct bacterial populations were used, one phenotypically normal, the other composed of persister cells.  
Persister cells are defined as slower growing and are less susceptible to antibiotics \cite{Mikkaichi2019}.
Mikkaichi's model analyzed the interaction of vancomycin, host immune function and differences in bacterial phenotypic behavior to identify the key factors that lead to persistent bacteremia.  Using this ensemble modeling approach, Mikkaichi and colleagues identified decreased immune clearance rate of persister cells as the primary determinant of persistent bacteremia \cite{Mikkaichi2019}. 

Vera-Yunca et al. developed a two-state semi-mechanistic PKPD model comprising a growing,
drug-susceptible population and a non-growing, drug-insusceptible dormant state. The model was built on in vitro time–kill
experiments with linezolid and vancomycin against MRSA and subsequently adjusted to data from a neutropenic mouse pneumonia model.
The in vivo drug-effect parameters were obtained by re-estimating scaling
ratios relative to the in vitro estimates rather than independently.
Whereas Mikkaichi et al. incorporate host immune clearance explicitly, the neutropenic design of Vera-Yunca's model excluded it.


The present study uses a dual-compartment mathematical Pharmacokinetic/pharmacodynamic model (PKPD)
to examine the interplay of the factors that lead to resolution of bacteremia in the setting of osteomyelitis and bacterial endocarditis.
This type of model is used to connect drug exposure to drug response.
Linking time course of drug concentration to the time course of drug effect allows the model to simulate what happens during treatment.
Two MRSA populations are used across the two compartments.  One MRSA subpopulation is vancomycin resistant
while the other is sensitive to vancomycin.  The host system is characterized by a clinical diagnosis of bacterial
endocarditis with the initial bacteremia secondary to osteomyelitis of the right ankle.
The model evaluates how these two phenotypically distant MRSA sub populations grow in the environment of antibiotic treatment,
and immune response of in-silico patients. Using well understood population averages of clinical relapse in a specific clinical scenario, the techniques of Monte carlo sampling and non-linear dynamics are used to delineate what properties of the drug–host–reservoir system produce durable bacteremia clearance or post-treatment relapse.  By integrating bacterial dynamics, antibiotic pharmacology, and the host innate immune response within a single modeling framework, this approach provides a means of studying the interplay between these forces in a way that human clinical trials cannot.  
 
\maketitle
\section{Methods}

\subsection{Case Study: Clinical Scenario}
Bacterial endocarditis, an infection of the heart valves, frequently develops as a complication of a deep, unresolved bone infection known as osteomyelitis. In this scenario, bacteria from the infected bone continuously leak into the bloodstream creating a bacteremia. The bacteremia allows bacteria to travel to the heart and attach to its valve structures. The following case illustrates this clinical pathway and serves as the foundation for the proposed pharmacokinetic and pharmacodynamic simulation study.

The case involves a 62-year-old man with type 2 diabetes who was admitted to the hospital after outpatient antibiotic treatment failed to clear a chronic, non-healing wound on his right ankle. Because diabetes impairs both circulation and immune function, the infection had spread deep into the bone, causing right ankle osteomyelitis. Upon admission, the patient was ill, presenting with progressive fatigue and high, intermittent fevers peaking at 39°C. Physical examination showed a wound on the right ankle draining clear blood, while laboratory work revealed elevated C-reactive protein (a marker of severe systemic inflammation) and a high white blood cell count. Although a physical examination of his heart was unremarkable, a transesophageal echocardiogram (TEE), a specialized ultrasound probe passed down the esophagus to view the heart up close—showed thickening of the mitral valve’s posterior leaflet, confirming the diagnosis of bacterial endocarditis.

Based on a presumptive diagnosis of staphylococcal endocarditis, blood cultures were drawn, and the patient was started on intravenous vancomycin at 1,000 mg every 8 hours. Despite this treatment, he continued to spike fevers. On the third day of hospitalization, blood culture results identified the organism as methicillin-resistant *Staphylococcus aureus* (MRSA) with a vancomycin Minimum Inhibitory Concentration (MIC) of 32 mg/L. This high MIC indicated that the strain was resistant to vancomycin, explaining why the patient was not improving. Laboratory testing showed that the bacteria were sensitive to linezolid. Vancomycin was discontinued, and the patient was switched to linezolid at 600 mg every 12 hours. Over the next 48 hours, his fever curve normalized.

To complete the treatment for MRSA endocarditis, the patient required a 6-week course of linezolid in accordance with Infectious Diseases Society of America (IDSA) guidelines. The primary clinical endpoint of this therapy is the complete resolution of bacteremia through the eradication of the primary infection source, which in this patient is the right ankle osteomyelitis. However, even with effective antibiotic therapy, 20\% to 30\% of patients experience a relapse of bacteremia because the infected bone continues to seed bacteria into the blood unless surgical intervention is performed to remove the source. The proposed dual-compartment pharmacokinetic/pharmacodynamic model simulates this exact dynamic.  It tracks how linezolid distributes through body tissues (pharmacokinetics) and suppresses bacterial growth (pharmacodynamics) in the presence of a persistent, seeding bone infection.

\subsection{Bacterial System and ODE Model}
The dynamical framework consists of four state variables, $S$, $R$, $S^r$, and $R^r$, each representing the four distinct
subpopulation of MRSA in the two compartments.  Two of the subpopulation reside in the blood, vancomycin sensitive, $S$ and
resistant MRSA, $R$ and the other two represent vancomycin sensitive and resistant
MRSA in the reservoir, $S^r$ and $R^r$.  All four subpopulations are governed by a coupled system of four ordinary differential equations.
Equations~\eqref{eq:dSdt} and \eqref{eq:dRdt} relate to processes in the vascular compartment, Equations~\eqref{eq:dSrdt} and \eqref{eq:dRrdt} 
relate to processes in the reservoir compartment.  The natural death rate of bacteria was omitted from the model due to it's small size
compared to the other variables.


\begin{equation}
\label{eq:dSdt}
\frac{dS}{dt}
= \rho_{s}\, S \left( 1 - \frac{S + R}{K} \right)
- k_{\text{imm}}\, S
- E^{V}(C_V)\, S
- E^{L}(C_L)\, S
+ k_{r\to b}\, S^r .
\end{equation}


\begin{equation}
\frac{dR}{dt}
= \rho_{r}\,R \left( 1 - \frac{S + R}{K} \right)
- k_{\text{imm}}\, R
- E^{L}(C_L)\, R
+ k_{r\to b}\, R^r .
\label{eq:dRdt}
\end{equation}



\begin{equation}
\frac{dS^r}{dt}
= \rho_{s}^{r} S^{r} \left( 1 - \frac{S^r + R^r}{K^r} \right)
- 0.5\, k_{\text{imm}}\, S^r
- E^{V}(\phi_{V} C_V)\, S^r
- E^{L,r}(\phi_{L} C_L)\, S^r
- k_{r\to b}\, S^r .
\label{eq:dSrdt}
\end{equation}


\begin{equation}
\frac{dR^r}{dt}
= \rho_{r}^{r} R^{r} \left( 1 - \frac{S^r + R^r}{K^r} \right)
- 0.5\, k_{\text{imm}}\, R^r
- E^{L,r}(\phi_{L} C_L)\, R^r
- k_{r\to b}\, R^r .
\label{eq:dRrdt}
\end{equation}



Each drug effect term is a sigmoidal $E_{max}$ relationship with a Hill
coefficient of unity, evaluated at the free-drug concentration reaching
that compartment:
\begin{equation}
E^{i}(C) = \frac{E_{\max}^{i}\, C}{EC_{50}^{i} + C},
\qquad i \in \{V, L\},
\label{eq:hill}
\end{equation}


where $C_V$ and $C_L$ are the free vancomycin and linezolid
concentrations in blood from Eq.~\eqref{eq:pk}, and the reservoir sees
only the penetrated fractions $\phi_V C_V$ and $\phi_L C_L$. The maximum rate of linezolid killing in the
reservoir is reduced by an independent efficacy factor,
$E_{\max}^{L,r} = \alpha\,E_{\max}^{L}$ with $\alpha = 0.45$.   This term reflects the
reduced activity of linezolid against slow-growing, biofilm-embedded and
small-colony-variant \textit{S.~aureus} in bone. The factor $\alpha$ is seperate
from the reservoir growth rates and is an assumed value, estimated by calibration.
It gives an effective reservoir maximum linezolid killing rate of
$0.45 \times 1.97 \approx 0.89~\mathrm{h}^{-1}$. Vera-Yunca et al.\ estimated that linezolid
$E_{\max}$ in a neutropenic mouse pneumonia model was 0.303 times its in vitro value (95\% CI
0.269--0.330), an in vivo maximum of approximately $0.60~\mathrm{h}^{-1}$.   and attributed the
reduction to factors other than drug exposure at the infection site \cite{VeraYunca2025}.
Because the present model represents reduced bone exposure explicitly through the penetration
fraction $\phi_L$, part of that attenuation is already accounted for, and $\alpha$ is
correspondingly larger than 0.303; it lies within the 0.2--0.6 range over which that
attenuation is uncertain.
$E^{L,r}$ denotes the resulting effect term.

Both compartments have distinct logistic growth dynamics and the carrying
capacities of each compartment is shared between the sensitive and resistant species.
The logistic term depends on the combined density $S+R$ (blood) or $S^r+R^r$ (reservoir), so the strains compete directly for the same niche.
Vancomycin acts on the sensitive strain only.
The immune response in the reservoir is scaled to a fraction of the
blood compartment. As detailed previously, the reservoir acts as a seeding
source for both sensitive and resistant populations entering the blood
compartment. Return movement from blood to reservoir is neglected.



\subsection{Model Calibration}
The initial condition of the reservoir assumes inoculum of $100$ for sensitive and resistant bacterial species respectively.
The blood compartment at the beginning of the simulation
is assumed to be free of pathological bacterial species $\mathit{S(0) = 0}$ and $\mathit{R(0) = 0}$, which is consistent
with normal human physiology.  Values for vancomycin and linezolid $EC_{50}$ and $E_{max}$ reported in the literature
are variable and likely due to methodological differences between studies.  This limitation makes it more challenging
to use all of the limited amount of data currently available.

To provide internal consistency, growth rates and pharmacodynamic values ($E_{max}$ and $EC_{50}$)
for both vancomycin and linezolid were taken from Vera-Yunca and colleagues \cite{VeraYunca2025}.
The study used in vitro time-kill and in vivo neutropenic mouse lung data to fit growth rates
and antibiotic PD parameters.  The in-vitro data was fit to the mouse data with a scaling factor.
In our study we used the in-vitro data to represent the blood compartment and the in-vivo, the reservoir compartment.
We used a scaling factor to adjust the growth rate of vancomycin sensitive MRSA in the reservoir and linezolid's $E_{max,L}$.
No changes were made to the $EC_{50,L}$.  The scaling factor used was based on calibration and not the same value used by
Vera-Yunca and colleagues.  The percentage of simulated cases where bacteremia is present by the end of the followup period
are based on population averages \cite{garciadelpozo2018factors}.  

\subsection{Blood Compartment}
The blood compartment represents the circulatory system of an individual with blood moving rapidly between
the organs and throughout the body.  There is no attenuation of bacterial killing by the immune system.
Carrying capacity is shared between sensitive and resistant bacteria and represents the maximum bacterial load that can be sustained within
the vascular compartment. The carrying capacity is based on literature estimates from \cite{nieman2022bacterial}.
The growth rate of sensitive SA is set to 1.47 $h^{-1}$, the in vitro estimate for MRSA reported by
Vera-Yunca and colleagues \cite{VeraYunca2025}, treating well-mixed, nutrient-replete blood as in-vitro-like.
This value corresponds to a doubling time of 28 minutes.  The resistant SA growth rate is scaled by a 10\% fitness cost to 1.323 $h^{-1}$, doubling time of 31 minutes.
The microbiological parameters are detailed in \cref{tab:Comparsms}.


\begin{figure}
    \centering
    \includegraphics[width=1.0\linewidth]{graph.png}
    \caption{Two Compartment Model} 
    \label{fig:graph.png}
\end{figure}


\subsection{Reservoir Compartment}
The reservoir compartment models the osseous environment and it the initial source of infection after having spread from the ankle wound to the long bone.
Infected bone undergoes localized inflammation and creates sequestrum (necrotic bone surrounded by infected granulation tissue)
with poor immune cell access and restricted antibiotic diffusion \cite{MomoduSavaliya2023}. Staph aureus forms biofilms that adhere
to the sequestrum forming an organized structure with the bacteria surrounded by a matrix.  Biofilm can release planktonic microbes from their
surface and seed the circulation, eventually leading to new biofilms at other locations such as the mitral valve.
Staph aureus forms Small Colony variants (SVC) that can reside inside osteoblasts, providing further protection
from antibiotics and the immune system\cite{Kahl2016SCV}. SCVs are slow-growing phenotypic variants that can form colonies and have been isolated
from infections such as osteomyelitis and endocarditis \cite{Kahl2016SCV}.  These factors lead to slower growth and lower metabolic activity
in bone driving a state of chronic persistence.  The mathematical model takes these physiological and microbiological characteristics
into consideration by assuming lower growth rates between the two compartments. The growth rate for sensitive bacteria is 0.845 $h^{-1}$ in the
reservoir, corresponding to a doubling time of 49 minutes; the resistant strain carries the same 10\% fitness cost as in blood
(0.761 $h^{-1}$, 55 minutes).  The lack of data on bone growth rates requires the use of a surrgoate value to
represent bone.  The reservoir value is the in vivo maximum growth rate that Vera-Yunca and colleagues estimated
in a mouse lung model, 42\% lower than their in vitro rate of 1.47 $h^{-1}$ \cite{VeraYunca2025}. Its use for bone
assumes that the in vivo reduction in growth applies similarly to poorly perfused, biofilm-associated bone.
Exchange between compartments is assumed to be one-way and governed by a first-order rate constant $k_{r \to b}$ (reservoir to blood).
Return movement of bacteria from blood into the reservoir is negligible compared with shedding from the established bone focus into the blood.
This is an assumed value due to the lack of human experimental measurements.  Similarly to the blood compartment, the reservoir carry capacity is shared
between sensitive and resistant bacteria, however, unlike the blood compartment it is assumed to be higher.
This assumption is based on the assumed presence of greater nutrient sources in necrotic bone and values
reported in the literature \cite{VeraYunca2025}.





\begin{equation}
R^{b} + S^{b}
\,
\xleftarrow{\;k_{r \to b}\;}
\,
R^{r} + S^{r}
\end{equation}



\begin{table}[H]
\centering
\large %Increase the text size within this table
\renewcommand{\arraystretch}{1.35} % Adds vertical breathing room between rows
\caption{Compartmental Parameter Values}\label{tab:Comparsms}
\vspace{0.5cm}
\begin{tabular}{l l l l}
\hline
\textbf{Parameter} & \textbf{Units} & \textbf{Nominal Value} & \textbf{Reference} \\
\hline
$\rho_s^{b}$ & $h^{-1}$ & 1.47 & \cite{VeraYunca2025}\\
$\rho_r^{b}$ & $h^{-1}$ & $1.323$ &  10\% fitness cost  \\
$\rho_{s}^{r}$         & $\text{h}^{-1}$     & $0.845$             & \cite{VeraYunca2025} \\
$\rho_{r}^{r}$         & $\text{h}^{-1}$     & $0.761$             & 10\% fitness cost\\
$B_{\max}^{b}$ & CFU/mL & $6 \times 10^3$ & \cite{nieman2022bacterial} \\
$B_{\max}^{r}$         & CFU/mL              & $1 \times 10^{7}$   & \cite{VeraYunca2025}\\
$k_{r \to b}$          & $\text{h}^{-1}$     & $1 \times 10^{-7}$ & \cite{Kostakioti2013} \\
\hline
\end{tabular}
\end{table}

 

\subsection{Immune Response}
The onset of infection activates the body's innate immune system. White blood cells (WBCs), particularly
neutrophils, which are largely quiescent prior to infection, become activated and migrate to sites of bacterial invasion.
Once at the site of infection, these cells release bactericidal substances, including reactive oxygen species and proteolytic enzymes,
before undergoing cell death. Pro-inflammatory cytokines and chemokines released at the site of infection stimulate
the bone marrow to increase WBC production and mobilization, leading to rapid delivery of immune cells through the circulatory
system to infected tissues.

Host immune clearance is represented as a constant, first-order killing
term. All four bacterial subpopulations in both blood and
reservoir are cleared at a rate proportional to their own density and to
a single rate constant, $k_{immune}$ (baseline $0.18~\mathrm{h^{-1}}$).
The reservoir immune response is scaled to 0.5 of the response in blood. This
reduction models biofilm exclusion of immune cells and the intracellular protected state of MRSA.
The model for immune response is an assumed and simple representation
of a highly complex biological process that depends on a multitude of factors.


\subsection{Antibiotic Pharmacokinetics}
A standard protocol for the treatment of infectious endocarditis is used, based on the recommendations
of the Infectious Disease Society of America (\cite{Minter2023}).
Vancomycin is the first line agent and linezolid secondarily once vancomycin resistance was been recognized.
Vancomycin, a weakly bactericidal agent, inhibits bacterial cell wall synthesis and linezolid, a bacteriostatic agent,
inhibits protein synthesis. Studies have shown that linezolid possesses both intracellular and extracellular activity, against Staphylococcus\cite{sandberg2010intra}.
During simulation, guideline doses, $\mathrm{Dose}_{i}$, and dosage time intervals, $\tau_{i}$, for vancomycin and linezolid are used \cite{vancopk_review}.  Total drug administration time is parameterized by Duration. 

Serum drug concentrations are simulated with a one-compartment PK model for both Vancomycin and Linezolid.
Vancomycin is approximately 50\% protein-bound in plasma, whereas linezolid is more modestly protein-bound ($\sim$30\%, i.e., $\sim$70\% free fraction), a difference the model accounts for via each drug's free-fraction parameter ($f_u$ = 0.50 for vancomycin, 0.70 for linezolid) when converting total plasma concentration to $C_{i}$, the active free-drug concentration.  Antibiotic clearance, $CL_{i}$ is consistent with normal kidney function. Population averages are used for the apparent volume (liters) of distribution, $V_{d}$. 


Between doses, $\mathrm{Dose}_{i}$, the concentration decays via first-order kinetics with elimination rate $k_{e,i} = CL_i / V_{d,i}$. At each dosing time $t_k = k\,\tau_i$ the concentration steps up by $\mathrm{Dose}_i / V_{d,i}$ and then decays again,
so that within the $k$-th dosing interval


\begin{equation}
C_i(t) = C_i(t_k^{+})\, e^{-k_{e,i}(t - t_k)},
\qquad t_k \le t < t_{k+1},
\label{eq:pk}
\end{equation}
with the post-dose value set by the accumulation rule
\begin{equation}
C_i(t_k^{+}) = C_i(t_k^{-}) + \frac{\mathrm{Dose}_i}{V_{d,i}},
\qquad C_i(0^{-}) = 0 ,
\label{eq:pkdose}
\end{equation}


for $i \in \{V, L\}$ (vancomycin, linezolid). Each drug is dosed only within
its own course window. Outside the window no further doses are added and the
concentration decays toward zero. Multiplying by the unbound fraction
$f_{u,i}$ gives the free-drug concentration that drives pharmacological effect.



\begin{table}[H]
\centering
\large % Increase text size within the table
\renewcommand{\arraystretch}{1.35} % Adds vertical breathing room between rows
\caption{Table of Pharmacokinetic Parameter Values}\label{tab:PKparams}
\vspace{0.3cm} % Adds space between caption and table
\begin{tabular}{l l l l} % Change to 'c' if you want centered column contents
\hline
\textbf{Parameter} & \textbf{Units} & \textbf{Nominal Value} & \textbf{Reference} \\
\hline
$E_{\max}^{V}$ & $\mathrm{h^{-1}}$ & 1.7 & \cite{VeraYunca2025}\\
$E_{\max}^{L}$ & $\mathrm{h^{-1}}$ & 1.97 & \cite{VeraYunca2025}\\
$EC_{50}^{V}$ & mg/L & 0.25 & \cite{VeraYunca2025}\\
$EC_{50}^{L}$ & mg/L & 0.56 & \cite{VeraYunca2025}\\
$V_d^{V}$ & Liters & 50 & \cite{Matzke1986} \\
$V_d^{L}$ & Liters & 45  &  \cite{dryden2011linezolid} \\
$CL_{V}$ & ml/min & 83 (5.0 L/h) & \cite{Matzke1986} \\
$CL_{L}$ & ml/min & 100 (6.0 L/h) & \cite{dryden2011linezolid} \\
$\phi_{LZD}$           & --                  & $0.50$              & \cite{dryden2011linezolid} \cite{kutschalissberg2003linezolid} \\
$\phi_{VAN}$           & --                  & $0.15$              & \cite{MomoduSavaliya2023} \cite{spellberg2012systemic} \\
$\alpha$               & --                  & $0.45$              & \cite{VeraYunca2025} \\
\hline
\end{tabular}
\end{table}



There is significant variability reported in the literature on the degree of bone penetration of linezolid,
although clinicians generally view it's bone penetration as favorable and significantly greater than vancomycin.
Linezolid bone penetration was set to 50\% creating free linezolid concentration from 1.2 to 5.9 mg/L, averaging 2.9 mg/L.
This estimate compares favorably
to data reported by \cite{kutschalissberg2003linezolid} using the same dosing regimen.  That study reported a mean bone
concentration of 3.9 mg/L (range: 1.2 - 8.3 mg/L).
Vancomycin bone penetration was set to 15\%.


During simulation the steady state peak Vancomycin concentration reached 36.3~mg/L and a trough concentration of 16.3~mg/L.  The total-drug
area under the concentration-time curve over 24 hours of
$\mathrm{AUC}_{24,\mathrm{ss}} = 600~\mathrm{mg\cdot h/L}$,
independent of volume of distribution or half-life. This falls at the upper
edge of the standard clinical target of 400--600~mg$\cdot$h/L for MRSA
bacteremia and endocarditis\cite{rybak2020therapeutic}. Accounting for the drug's free fraction
($f_u = 0.50$), the corresponding free-drug exposure driving bacterial kill
in the model is $\mathrm{AUC}_{24,\mathrm{ss}}^{\text{free}} = 300~\mathrm{mg\cdot h/L}$.
Linezolid steady-state total-drug trough concentration was calculated at 3.4~mg/L,
which is within the recommended range of 2 to 8 mg/L if the MIC is less than 2 mg/L\cite{alsultan2019determining}. 


\subsection{Definition of treatment Failure and Success}

Treatment failure is defined as a bacteremia state at the end of the simulation and success as resolution of the bacteremia. 
Laboratory testing is unable to detect the presence of bacteria below a pre-defined level.
The model uses the Level of Detection (LOD) to represent this fact.  When the concentration of MRSA drops below this level the model sets it's value to zero.  The limit of detection (LOD) was set at $10~\mathrm{CFU/mL}$ and implemented as an extinction boundary where any bacterial subpopulation that fell below the LOD are set to zero. MRSA cannot grow back from a fractional leftover remnant. 

An extinct blood subpopulation can be re-established if continued transfer from a detectable reservoir accumulated to the LOD. Treatment failure is evaluated by two endpoints, persistent
bacteremia and a relapsing bacteremia.  Persistent bacteremia is defined as detectable levels of MRSA in the blood at the end of the linezolid course and a relapsing bacteremia is defined as a the absence of MRSA in the blood at the end of treatment, but a return to detectable levels of MRSA at the end of the simulation.  This results is due to a reseeding of the blood compartment from the reservoir. Let $R_B(t)$ denote the MRSA
density in blood at time $t$ (days) and $t_L$ the time of linezolid completion.
A persistent bacteremia is defined as detectable resistant MRSA in blood at
the end of linezolid treatment,
\begin{equation}
    \mathcal{P} = \Big\{ MRSA_B(t_L) \ge \mathrm{LOD} \Big\}.
    \label{eq:persistence}
\end{equation}

Relapse is defined as clearance of resistant MRSA from
blood at linezolid completion followed by a return to detectable levels
during the 14-day follow-up,
\begin{equation}
    \mathcal{R} = \Big\{ MRSA_B(t_L) < \mathrm{LOD} \Big\} \cap
    \Big\{ \exists\, t \in (t_L,\, t_L + 14] : MRSA_B(t) \ge \mathrm{LOD} \Big\}.
    \label{eq:relapse}
\end{equation}

These two conditions at $t_L$, are mutually exclusive, and simulation results are reported separately for each outcome.  Separate Monte Carlo simulations are performed for each parameter over the experimental period of 81 days. This produces a spread of possible outcomes (trajectories) for each of the four groups of MRSA.


\subsection{Threshold localization with Brent's method}

Resolution of bacteremia or no resolution of bacteremia were two clinical outcomes used to calculate the threshold value.  The $EC_{50,L}$, $E_{\max,L}$, and $k_{immune}$ that determine the boundary between the two outcomes were
located using the Brent algorithm. Let $\theta$ denote the
parameter being varied and $B_{blood}(t_f;\theta)$ the bacteremia 
burden at the final simulated time $t_f$. One complete 81-day ODE simulation
was run for each evaluation of
\begin{equation}
    F(\theta) = R_{res}(t_f;\theta) - L,
    \qquad L = 10~\mathrm{CFU/mL},
    \label{eq:brent-objective}
\end{equation}
so that $F(\theta) < 0$ denotes bacteremia resolution and $F(\theta) > 0$ denotes
the no resolution of bacteremia. Grids spanning
$[0.05,\, 3.3]~\mathrm{mg/L}$ for $EC_{50,L}$,
$[0.70,\, 2.2]~\mathrm{h}^{-1}$ for $E_{\max,L}$, and
$[0.02,\, 0.65]~\mathrm{h}^{-1}$ for $k_{immune}$ were initially used to verify a
single monotone change of outcome. Adjacent grid points falling on opposite sides of the transition
supplied a sign-changing bracket $[a,b]$ with
\begin{equation}
    F(a)\,F(b) < 0.
    \label{eq:brent-bracket}
\end{equation}
SciPy's bracket-preserving \texttt{scipy.optimize.brentq} implementation was
then applied to $F(\theta)$ \cite{2020SciPy-NMeth}. The method combines bisection,
secant interpolation, and inverse quadratic interpolation, accepting an
interpolated step only when it remains within the current sign-changing
bracket and taking a bisection step otherwise. Each accepted evaluation
replaces the endpoint of matching sign, so the boundary remains enclosed
throughout the search. An absolute parameter tolerance of $10^{-3}$ in the
corresponding parameter units was used, and transition values are reported to
two significant figures, except for $E_{max,L}$, reported to three because its narrow
sampling distribution makes the difference between 1.85 and 1.9 material.

Brent's root-finding method can be used for both continuous and discontinuous functions \cite{brent1973}.  The
hard LOD rule that was established makes Eq.~\eqref{eq:brent-objective}
discontinuous. Brent \cite{Brent1971} does not require the function to be continuous
over the bracketing interval. For discontinuous differential systems, an event can first be
detected within an integration step and subsequently localized using
a bracketed root-finding procedure. Brent-type methods combine
interpolation with bisection to refine the event location, while
retaining a bracket around the event \cite{Brent1971,WolframEventDetection}. When the function is discontinuous, a
sign-changing interval may contain no true zero; consequently, the
algorithm instead seeks a sufficiently small interval in which the
function takes both non-negative and non-positive values.Any trajectory crossing below the LOD is reset to zero and
cannot regrow autonomously.  As a consequence of the LOD the terminal burden can jump from zero to a large concentration of bacteria (CFU/mL) without attaining the exact LOD. ($10~\mathrm{CFU/mL}$). For
$EC_{50,L}$, simulations immediately below and above the located boundary
produced terminal burdens $(R_b, R_{res})$ of $(0,0)$ and approximately
$(5.2 \times 10^{3},\, 8.8 \times 10^{6})~\mathrm{CFU/mL}$, respectively,
confirming that $F$ does not pass continuously through zero. For each of the threshold evaluations the returned value represents the boundary where the sign of the function changes.  It is not interpreted as a parameter satisfying $F(\theta) = 0$.
Bacteremia lies above this boundary for $EC_{50,L}$ and below it for $E_{\max,L}$ and $k_{immune}$.

\subsection{Probability of Bacteremia Clearance }
In the prior section the Brent method is used to calculate the boundary point that separates the two clinical outcomes.  The probability of each parameter is modeled with a log normal distribution, which clusters most values around the mean with a small number of parameter values spread far above in the tail \cite{limpert2001lognormal}. For the log-normal $EC_{50,L}$, $E_{max,L}$, and $k_{immune}$ distributions, arithmetic mean $m$ and coefficient of variation $CV$ are converted to log-space
parameters as
\begin{equation}
\sigma_{\ln}=\sqrt{\ln(1+CV^2)}, \qquad
\mu_{\ln}=\ln(m)-\frac{\sigma_{\ln}^{2}}{2}.
\label{eq:lognormal-parameters}
\end{equation}

The distance between the boundary and the specific parameter value is calculated in terms of case to case variation. In the example of drug potency, $EC_{50,L}$ the bacteremia does not resolve by the end of the simulation if the $EC_{50,L}$ is located above the boundary. For the other two parameters, maximum drug effect and immune response the bacteremia is not resolved by the termination of the simulation if the values lie below the boundary.  
Each boundary point $T$, is standardized log-space distance from the
specific parameter distribution is calculated by
\begin{equation}
z_T=\frac{\ln(T)-\mu_{\ln}}{\sigma_{\ln}}.
\label{eq:lognormal-switch-z}
\end{equation}

The symbol $\Phi$ denotes the standard-normal cumulative distribution function and represents the
probability that bacteremia does not resolve by the end of the simulation is
\begin{equation}
p_{bacteremia}=1-\Phi(z_T) \quad \text{for } EC_{50,L}, \qquad
p_{bacteremia}=\Phi(z_T) \quad \text{for } E_{max,L}\text{ and }k_{immune},
\label{eq:lognormal-bacteremia-tail}
\end{equation}

Changes in the distribution used to model the parameter value will alter the distance and location of the parameter in relation to the boundary point.  Change in the location of the probability mass with respect to the boundary will affect whether the bacteremia has resolved by the end of the simulation.




\maketitle
\section{Results}

\subsection{Pre-Treatment Period}
Prior to initiation of vancomycin, bacterial populations developed in the reservoir and seeded the blood compartment. The vancomycin-sensitive \textit{S. aureus} strain in blood had growth rate $\rho_s = 1.47\text{ h}^{-1}$ while the resistant strain had $\rho_r = 1.323\text{ h}^{-1}$, a 10\% fitness cost. In the bone reservoir, replication rates were reduced to $\rho_s^r = 0.845\text{ h}^{-1}$ for the sensitive strain and $\rho_r^r = 0.761\text{ h}^{-1}$ for the resistant strain, the same 10\% cost.
At the start of vancomycin the baseline reservoir held about $8\times10^{6}$~CFU/mL of the
sensitive strain and $3\times10^{4}$~CFU/mL of the resistant strain.
Once vancomycin has eradicated the senstive strain the resistant population grows up to the carrying capacity of the compartment.

The simulation represents progression from osteomyelitis, with associated wound of the ankle, to secondary bloodstream infection and subsequent development of bacterial endocarditis. Vancomycin administration (1000~mg every 8 hours) models clinically recommended steady-state peak and trough plasma concentrations and drives rapid killing of susceptible \textit{S. aureus} subpopulations ($S_b$ and $S_{\text{res}}$). Vancomycin exerted no growth inhibition on resistant strains ($R_b$ and $R_{\text{res}}$).

\subsection{Vancomycin Treatment Period}

The 4-day vancomycin course (\texttt{van\_duration} $= 96$ h, starting at \texttt{vanco\_start} $=
504$ h) drives the vancomycin-sensitive strain to extinction in
both blood ($S_b$) and reservoir ($S_{res}$) compartments before
the initiation of linezolid.
Following laboratory confirmation of vancomycin-resistant MRSA, therapy was switched to Linezolid (600~mg every 12 hours),
yielding recommended steady-state serum concentrations. Linezolid administration drove a reduction in resistant bacterial
counts across both compartments. However, clearance
from the reservoir lagged significantly behind vascular clearance.


\subsection{Linezolid Potency}
To assess how variation in linezolid potency affects predicted
\textit{Staph aureus} bacterial kinetics, $EC_{50,L}$ was sampled from
a log-normal distribution with mean $0.56~\mathrm{mg/L}$ and coefficient
of variation $\mathrm{CV} = 0.60$ \cite{limpert2001lognormal}. The mean was taken from the linezolid pharmacodynamic
estimate of Vera-Yunca et al.\ ($0.56~\mathrm{mg/L}$) \cite{VeraYunca2025}. All other parameters were held at baseline. Each of the $n = 1000$
independently sampled $EC_{50,L}$ values were propagated through the model from blood-sterile, reservoir-seeded
initial conditions ($S_b(0){=}R_b(0){=}0$, $S_{res}(0){=}R_{res}(0){=}100$
CFU/mL) to create a distribution of simulated trajectories \cref{fig:ec50_kinetics}.
Sensitive SA in both compartments is eradicated by vancomycin by the end of it's treatment course.
Many of the simulations the vancomycin resistant SA are eradicated in the blood compartment, but persist, albeit at a lower level,
in the reservoir compartment.  Termination of linezolid allows the reservoir to re-seed the blood compartment causing the relapse
of the bacteremia.

\begin{figure}[H]
    \centering
    \includegraphics[width=1.0\linewidth]{figures/mc_ec50_lzd_sweep_kinetics.png}
    \caption{Monte Carlo sweep of $EC_{50, L}$ ($n = 1000$): full
    time-course kinetics for $S_b$, $R_b$, $S_{res}$, and $R_{res}$, shown
    as individual trajectories (light), median (solid), and 5th--95th
    percentile envelope (shaded).}
    \label{fig:ec50_kinetics}
\end{figure}

Increasing $EC_{50,L}$ represents decreasing linezolid potency and a greater probability of a lack of bacteremia resolution.  Bacteremia relapse occurred in 22\% of simulated cases, persistent bacteremia in 0.2\% (two patients), and complete resolution in 78\%. The $EC_{50,L}$ value that is the transition point between clinical resolution from failure is reported as $EC_{50,L}=0.76~\mathrm{mg/L}$ \cref{fig:ec50_switch_population}.  The upper panel shows the bacterial load for each compartment at the termination of each simulation, given the value of $EC_{50,L}$. The lower panel illustrates the distribution of $EC_{50,L}$ values from the given distribution.  As illustrated in \cref{fig:ec50_switch_population} the majority of $EC_{50,L}$ probability mass is localized to the side of the transition that favors clinical resolution.


\begin{figure}[H]
    \centering
    \includegraphics[width=1.0\linewidth]{figures/mc_ec50_lzd_switch.png}
    \caption{Upper panel: Final bacterial burden in each compartment, Lower panel: Distribution of $EC_{50,L}$ values color coded with final clinical outcome.}
    \label{fig:ec50_switch_population}
\end{figure}
The methodology for the calculation of the transition point is detained in section D.





\subsection{Linezolid Maximum Effect}

To assess how variation in linezolid maximum killing capacity affects
predicted \textit{Staph aureus} bacterial kinetics, $E_{max,L}$ was sampled
from a mean-corrected log-normal distribution with arithmetic mean
$1.97~\mathrm{h}^{-1}$ and coefficient of variation $\mathrm{CV} = 0.080$.
Published $E_{max,L}$ linezolid estimates span two distinct ranges: $0.735$--$0.754~\mathrm{h}^{-1}$ in Sandberg et al.\ \cite{sandberg2010intra} and $1.77$--$2.22~\mathrm{h}^{-1}$ in Scheerans et al.\ \cite{Scheerans2014}. The present analysis uses the in vitro estimate of Vera-Yunca et al.\ ($E_{max,L} = 1.97~\mathrm{h}^{-1}$, $EC_{50,L} = 0.56~\mathrm{mg/L}$) \cite{VeraYunca2025}. Each of $n = 1000$ independently sampled $E_{max,L}$ values was propagated through the model over the complete 81-day treatment and
follow-up course \cref{fig:emax_kinetics}.  

\begin{figure}[H]
    \centering
    \includegraphics[width=1.0\linewidth]{figures/mc_emax_lzd_sweep_kinetics.png}
    \caption{Monte Carlo sweep of $E_{\mathrm{max}, L}$: $S_b$, $R_b$, $S_{\mathrm{res}}$ and $R_{\mathrm{res}}$, with individual relapsed trajectories (light), the relapsed-patient median (solid), and relapsed-patient 5th--95th percentile envelope (shaded).}
    \label{fig:emax_kinetics}
\end{figure}

Lower $E_{max,L}$ values represent weaker antibiotic effect and shift the outcomes to clinical relapse and persistence.  Bacteremia resolution occurred in 77\% of simulated patients, relapse in 23\% and no persistent bacteremia \cref{fig:emax_switch_population}.  The transition point between treatment response and failure was reported at $1.85~\mathrm{h}^{-1}$ as illustrated by the upper panel of \cref{fig:emax_switch_population}.  The lower panel shows the majority of $E_{max,L}$ probability mass is located above the $1.85~\mathrm{h}^{-1}$ transition point, favoring clinical resolution.   

\begin{figure}[H]
    \centering
    \includegraphics[width=1.0\linewidth]{figures/mc_emax_lzd_switch.png}
    \caption{Upper panel: Final bacterial burden in each compartment, Lower panel: Distribution of $E_{50,L}$ values color coded with final clinical outcome.}
    \label{fig:emax_switch_population}
\end{figure}

Methodology for transition point calculation is described in section D.




\subsection{Host Immune Response}

The immune response was modeled with a constant linear function ($k_{immune}$) with a value of $0.18~\mathrm{h}^{-1}$, which served as the mean of a log-normal distribution.  A coefficient of
variation $\mathrm{CV} = 0.80$ was used to provide a measure of biological variability.  Although unlikely to be a true representation of a highly complex system, it does provide a simple method to incorporate the innate immune response into this model. The value of $k_{immune}$ was arrived at by using the population average outcome for this specific clinical scenario and back calculating what value would be needed to obtain the expected clinical outcome based on population averages of 20-30\%.  This value is a calibration result and is not an independent prediction.  All remaining parameters were held at baseline. For each of
$n = 1000$ independently sampled $k_{immune}$ values were propagated through the
dual-reservoir model over the complete 81-day treatment and
follow-up course.  A distribution of trajectories was created for each of the four
bacterial subpopulations \cref{fig:immune_kinetics}.

\begin{figure}[H]
    \centering
    \includegraphics[width=1.0\linewidth]{figures/mc_immune_response.png}
    \caption{Monte Carlo sweep of $k_{immune}$ ($n = 1000$): full
    time-course kinetics for $S_b$, $R_b$, $S_{res}$, and $R_{res}$, shown
    as individual trajectories (light), median (solid), and 5th--95th
    percentile envelope (shaded).}
    \label{fig:immune_kinetics}
\end{figure}

Higher values of $k_{immune}$ bias the outcomes toward clinical resolution of the bacteremia.  Bacteremia resolution occurred in 72\% of simulated patients, relapse in 28\% and no persistent bacteremia.  The transition point between clinical resolution from failure is reported as $k_{immune}=0.093~\mathrm{h}^{-1}$ \cref{fig:immune_switch_population}.  The upper panel shows the bacterial load for each compartment at the termination of each simulation, given the value of $k_{immune}$. The lower panel illustrates the distribution of $k_{immune}$ values from the given distribution.


\begin{figure}[H]
    \centering
    \includegraphics[width=1.0\linewidth]{figures/mc_immune_response_switch.png}
    \caption{Upper panel: Final bacterial burden each compartment, Lower panel: Distribution of $k_{immune}$ values, color coded for final outcome.}
    \label{fig:immune_switch_population}
\end{figure}




%---------------------------------------------------------
\subsection{Clinical Response Determinants}

Each parameter $E_{max}$, $EC_{50,L}$ and $k_{immune}$ has a transition boundary separating bacteremia resolution from no resolution.
The $EC_{50,L}$ transition point is 0.76 mg/L, $E_{max}$ is $1.85~\mathrm{h}^{-1}$ and immune clearance is  $0.093~\mathrm{h}^{-1}$.
Individual Cases will vary around a mean value.  Each parameter is spread across 1000 virtual cases with a log-normal distribution.  The share
of patients on either side of the transition boundary will depend on how far the mean is from the specific boundary and the width of
the spread as determined by the coefficient of variation.  In this study the majority of individual cases land on the side where the bacteremia clears.

The share of patients that clear the bacteremia can be calculated directly.  The z-value is a measurement of how many standard deviations the boundary lies 
from the center of the distribution (on a log scale).
The normal distribution converts z into the expected share of patients that fall on one side of the transition point.

Equations~\eqref{eq:lognormal-switch-z}--
\eqref{eq:lognormal-bacteremia-tail}, evaluated at the 
transition points, gave
\begin{equation}
\begin{aligned}
z_{EC_{50}} &= 0.83,  & p_{persist} &= 1-\Phi(z_{EC_{50}})=20.2\%,\\
z_{immune} &= -0.59, & p_{persist} &= \Phi(z_{immune})=27.9\%,\\
z_{E_{max}} &= -0.74, &
p_{persist} &= \Phi(z_{E_{max}})=23.1\%.
\end{aligned}
\label{eq:evaluated-persistence-tails}
\end{equation}

For each parameter, the standardized log-space distance between the transition point and the centre of the
sampling distribution, $z_T = (\ln T - \mu_{\ln})/\sigma_{\ln}$, was converted through the standard normal
cumulative distribution function into the expected fraction of patients on the persistence side of the transition,
giving predicted fractions of 20.2\% for $EC_{50,L}$ ($z = 0.83$), 23.1\% for $E_{max,L}$ ($z = -0.74$), and
27.9\% for $k_{immune}$ ($z = -0.59$). These predictions agreed closely with the 1,000-patient Monte Carlo
cohorts, in which 222 $EC_{50,L}$ samples (22.2\%), 226 $E_{max,L}$ samples (22.6\%), and 278 $k_{immune}$
samples (27.8\%) fell on the persistence side. The largest discrepancy, 2.0 percentage points for $EC_{50,L}$,
is within the finite-sample variation expected for a cohort of this size (standard error about 1.3 percentage points).
These counts equal relapse plus persistent treatment failure ($220+2$, $226+0$, and $278+0$, respectively).
The transition point determines whether a patient fails, while the blood burden at linezolid completion determines
whether that failure is a relapse or persistent bacteremia.



\begin{table}[H]
    \centering
    \small
    \begin{tabular}{@{}lccccc@{}}
        \toprule
        Parameter & Mean & CV & Transition & Relapse & Persistent / Clinical Resolution \\
        \midrule
        $EC_{50,L}$ & 0.56 & 0.60 & 0.76 & 22.0\% & 0.2\% / 77.8\% \\
        $E_{max,L}$ & 1.97 & 0.080 & 1.85 & 22.6\% & 0.0\% / 77.4\% \\
        $k_{immune}$ & 0.18 & 0.80 & 0.093 & 27.8\% & 0.0\% / 72.2\% \\
        \bottomrule
    \end{tabular}
    \caption{Unified sampling inputs, reservoir-persistence
    transitions, and verified Monte Carlo outcomes. Units are mg/L for
    $EC_{50,L}$ and h$^{-1}$ for $E_{max,L}$ and $k_{immune}$.}
    \label{tab:current_summary}
\end{table}

Table~\ref{tab:current_summary} provides the cross-parameter comparison of
the calibrated inputs, persistence transitions, and simulated outcomes.

\Cref{fig:distribution_to_switch} visualizes the distributional calculation
behind these outcomes. Although the parameters have different units and
coefficients of variation, each calibration places approximately one-fifth to
three-tenths of its probability mass on the persistence side of the switch.
The theoretical tail probabilities and realized 1,000-patient fractions are
shown together so that finite-sample variation is visible.

\begin{figure}[H]
    \centering
    \resizebox{1.2\textwidth}{!}{\includegraphics{figures/distribution_to_switch.png}}
    \caption{Parameter distributions relative to their
    transition points. Histograms show the realized Monte
    Carlo samples, solid curves show the probability densities,
    and red shading identifies the persistence-side tail. The panels report
    the log-space $z$ value and corresponding normal-CDF tail probability for
    each mean-corrected log-normal distribution.}
    \label{fig:distribution_to_switch}
\end{figure}

The three-way clinical classification is compared directly in
\Cref{fig:stacked_outcomes}. Persistent treatment failure was essentially absent:
two $EC_{50,L}$ patients (0.2\%) still had detectable blood at linezolid completion,
and every other non-resolved trajectory first cleared blood during linezolid and
subsequently relapsed.

\begin{figure}[H]
    \centering
    \includegraphics[width=0.82\linewidth]{figures/stacked_outcomes.png}
    \caption{Stacked outcome proportions for the three 1,000-patient Monte
    Carlo analyses. Resolution, post-linezolid bacteremia relapse, and
    persistent treatment failure are mutually exclusive and sum to 100\%
    within each analysis.}
    \label{fig:stacked_outcomes}
\end{figure}

The cumulative incidence curves in \Cref{fig:relapse_time_curve} show that
the clinical relapses classified during the 14-day follow-up occur early in the post treatment phase of the simulation.
All three curves reach their final cohort relapse percentages within approximately two
days after linezolid completion. The denominator is the full 1,000-patient
cohort for each curve, not only patients who relapsed.

\begin{figure}[H]
    \centering
    \includegraphics[width=1.0\linewidth]{figures/relapse_time_curve.png}
    \caption{Cumulative incidence and timing of bacteremia relapse after
    linezolid completion. Each step marks the first post-treatment return of
    resistant blood bacteria to the $10~\mathrm{CFU/mL}$ LOD. Persistent
    treatment failures are excluded because their blood burden was already
    detectable when linezolid ended.}
    \label{fig:relapse_time_curve}
\end{figure}








\section{Discussion}

This study constructed a dual-compartment computational PKPD model to examine
how reservoir and blood populations respond to antimicrobial and immune
parameters. The current calibration targets a 20--30\% bacteremia
relapse band. Bacteremia relapse is defined in this context as the concentration of microbes to
be below the LOD at completion of linezolid treatment and to return to or above the LOD during follow-up.
Terminal reservoir persistence is considered as a separate endpoint. The $EC_{50,L}$ analysis
produced 220 relapses (22.0\%), 2 persistent treatment failures (0.2\%), and
778 resolutions (77.8\%), with a deterministic reservoir-persistence
transition point at $0.76~\mathrm{mg/L}$. The $E_{max,L}$ analysis produced 226
relapses (22.6\%), 0 persistent failures (0.0\%), and 774 resolutions
(77.4\%), with a transition point at $1.85~\mathrm{h}^{-1}$. The immune analysis
produced 278 relapses (27.8\%), 0 persistent failures (0.0\%), and 722 resolutions
(72.2\%), with a transition point at $0.093~\mathrm{h}^{-1}$.

The proximity of each distributional mean to its switch point influences what percentage of each cohort relapses.  The probability mass beyond
each switch point influences how many members of the cohort respond to Linezolid treatment. Moving
the $EC_{50,L}$ distribution upward, moving the $E_{max,L}$ or $k_{immune}$
distribution downward, or increasing variability so that more mass enters the
corresponding persistence tail would increase the fraction of the cohort that relapses.
Moving the distributions in the opposite directions would decrease the fraction of the cohort that relapses. This
relationship is not one-to-one with observed relapse because the trajectory
definition separates patients whose blood never clears during linezolid as
persistent treatment failures. In the present simulations, persistence-side
fractions were 22.2\%, 22.6\%, and 27.8\% for $EC_{50,L}$, $E_{max,L}$, and
$k_{immune}$, respectively; removal of 0.2\%, 0.0\%, and 0.0\% persistent
treatment failures yielded relapse rates of 22.0\%, 22.6\%, and 27.8\%.
The distribution equations predict persistence-side probabilities of 20.2\%,
23.1\%, and 27.9\%, respectively. Their differences from the realized
22.2\%, 22.6\%, and 27.8\% fractions are expected finite-sample variation in
cohorts of 1000 random draws.  The theoretical probabilities describe the
calibrated distributions, whereas the realized fractions describe this seeded
Monte Carlo cohort.
Accordingly, the calibrated mean, CV, distribution shape, switch location, and
clearance-before-reseeding requirement jointly determine the reported relapse
rate to antibiotic treatment.

This study provides a mathematical framework that allows one to better understand what variables affect how individual patients respond to antibiotic treatment
in the bacterial endocarditis setting. Despite significant progress in antibiotic development over
the past decades, successful treatment in the setting of osteomyelitis of the long bone remains challenging. 
Due to the technological challenges of collecting data from live patients in large numbers mathematical simulation using
what limited data are available may be the next best option.  Using what is understood regarding relapse rates based on large scale population studies
coupled with PKPD models, biofilms, and host immune responses, we can better understand the complex interactions between these factors and their impact on treatment outcomes. 
As our understanding of biofilm physiology, antimicrobial resistance and host immune responses improves, we can refine these models to better predict treatment outcomes. 


\section{Limitations and Future Directions}

This simulation study carries several limitations driven by current gaps
in empirical data.
The model relies on several assumptions concerning the bone reservoir
compartment. Key parameters such as inter-compartmental transfer rates,
intra-bone bacterial growth and death rates, and the baseline densities of
sensitive and resistant populations are not directly measured and were
estimated. Transfer between bone and blood was modeled as one-way, from
bone into blood only; re-seeding of the bone focus from the circulation was
neglected. While this simplification is biologically plausible (an
established focus sheds into the bloodstream far more than the circulation
seeds it), the reservoir-to-blood rate itself is an estimated value. None
of the values used in this part of the analysis were based on data.

Second, the host immune response is represented by a single first-order
clearance rate, $k_{immune}$, applied uniformly to each compartment (with
a lower effective rate in bone than in blood, reflecting reduced immune
access to a deep-seated focus). This rate had to be estimated, and the
model does not attempt to represent the underlying cellular or molecular
mechanisms of clearance. Neutrophil recruitment, killing dynamics,
the networks of cytokines\footnote{Signaling molecules that regulate
inflammation, cell growth, and immune cell activity.} and
chemokines\footnote{A family of small proteins that attract and guide
immune cells to specific sites of infection.} that coordinate them, are
important in guding the innate immune response and are omitted in favor of this single linear aggregate rate.
It is highly unlikely that immune function can be modeld by a linear function.

Third, the blood carrying capacity for \textit{S. aureus} is based on a
limited set of measurements using PCR-based quantification of bacterial
DNA, and the maximal carrying capacity used in the model was derived from
this same limited dataset.

Fourth, published estimates of $EC_{50}$ do not consistently distinguish
between total drug concentration and the free (unbound) fraction in
plasma; this model used free serum concentration throughout its
pharmacokinetic component, which may not match the basis of every $EC_{50}$
value drawn from the literature.

Fifth, the resistant strain carries a fitness cost in both compartments
--- 10\% in both blood ($\rho_r = 1.323$ versus
$\rho_s = 1.47~\mathrm{h}^{-1}$) and the reservoir
($\rho_r^r = 0.761$ versus $\rho_s^r = 0.845~\mathrm{h}^{-1}$) --- but
the cost is assumed rather than measured. The sensitive growth rates are taken
from a single study \cite{VeraYunca2025}: the in vitro rate for blood and the in vivo
mouse-lung rate for bone, neither of which was measured in bone.
The reservoir growth rates should therefore
be read as calibrated quantities chosen to reproduce the clinical
scenario, not as independently measured biological values. The reservoir
linezolid efficacy factor ($\alpha = 0.45$) is likewise a calibrated
quantity rather than a measured one, and the outcomes are sensitive to it:
small changes in $\alpha$ shift the relapse fraction of every analysis
substantially.

A further limitation is that the unified fixed baselines and Monte Carlo
means ($EC_{50,L}=0.56~\mathrm{mg/L}$, $E_{max,L}=1.97~\mathrm{h}^{-1}$, and
$k_{immune}=0.18~\mathrm{h}^{-1}$) have not been externally validated as a
joint parameter set. Any Monte Carlo percentages are conditional on this
parameterization and the specified distributions. The coefficients of
variation ($0.60$ for $EC_{50,L}$, $0.080$ for $E_{max,L}$, and $0.80$ for
$k_{immune}$) are outcome-calibrated model inputs, not independently estimated
population-variability parameters. External calibration
and validation against measured blood and bone burdens are required before
these outputs can be interpreted as patient risk estimates.

Finally, the in vitro bacterial doubling times used for calibration may
not fully reflect the metabolically dormant or small-colony-variant (SCV)
states typical of chronic biofilm infections, nor do they
necessarily reflect the actual replication rate of MRSA within the
circulatory system in vivo. All growth-rate calibration in this model is
ultimately anchored to in vitro measurements, and in vivo replication
dynamics -- particularly within a biofilm or vegetative focus -- remain a
source of uncertainty this model cannot resolve.

Additional limitations stem from constraints on available patient-level
data and the scope of the sensitivity analyses performed. Because
patient-level data on model parameters such as growth rates of
MRSA in blood pharmacodynamic and other physiological variables remain limited, the
model could not be formally validated against clinical outcome data. The
Monte Carlo sensitivity analyses were also univariate: each sweep varied
a single parameter ($EC_{50,L}$, $E_{max,L}$, or $k_{immune}$) while
holding every other parameter fixed at baseline. Real
inter-patient variability is likely joint and correlated.  Impaired renal function, can simultaneously reduce drug clearance and also alter immune status.  This univariate approach cannot capture
these interaction effects between parameters that vary together in an actual
patient population. Similarly, no attempt was made to introduce
inter-patient pharmacokinetic variability: volume of distribution,
clearance, and protein binding were held fixed at population-average
values for every simulated patient. Finally, the simulation
used a single fixed treatment protocol throughout, with no exploration
of missed doses, alternative dosing regimens, or source control (e.g. 
surgical debridement or hardware removal), all of which materially
affect outcomes in real bone and endocarditis infections.



\bibliographystyle{ieeetr}
\bibliography{references}






\end{document}
